% To be inserted as a robustness/convergent-evidence Results subsection

\subsection{External Convergent Validation}

The Koenig-Robert mouse-tracking data provided convergent evidence that $\rho$ tracks response competition outside the KH2017 semantic animal-categorization task. In the face/object study, face-like objects categorized as objects had higher stimulus-level $\rho$ than ordinary objects (ordinary $M=.571$, face-like $M=.602$, $\Delta=.031$, stimulus-cluster bootstrap CI [.012, .048], permutation $p=.002$). In the animal/object study, animal-like objects categorized as objects also had higher $\rho$ than ordinary objects (ordinary $M=.570$, animal-like $M=.724$, $\Delta=.154$, stimulus-cluster bootstrap CI [.103, .202], permutation $p<.001$). Pooled-object null simulations without a stimulus-level ambiguity predictor did not produce effects as large as the observed differences in either study.

The held-out path-likelihood comparison supported the intended interpretation of the model. The binary ambiguity-prior action model improved mean held-out NLL relative to the condition-only action model in both external studies (face/object: 46.003 vs. 46.342; animal/object: 51.041 vs. 51.574). The best raw NLL model was the trial-fitted action upper bound in both studies, underscoring that the external analysis evaluates whether $\rho$ tracks perceptual competition rather than whether the constrained action model is the most flexible raw trajectory curve-fitter.

SWOW-DE provided independent German association evidence for the original KH2017 item-level result. Direct item-to-category association coverage was sparse (11/19 items) and was not used as the main result. The prespecified main SWOW metric was therefore undirected network-proximity margin, which covered 18 of 19 items. Lower SWOW network margin predicted higher item-level $\rho$ ($r_s=-.590$, $p=.010$; Theil--Sen slope $=-.923$, bootstrap CI [-1.958, -.207]; permutation $p=.005$). The official SWOW PPMI--SVD embedding margin covered 11 items and showed the same direction ($r_s=-.536$, $p=.089$). The existing multisource transformer rank margin remained directionally consistent ($r_s=-.582$, $p=.009$).

The Schröder norm audit covered 13 of 19 KH2017 exemplars. No exact target--competitor typicality contrasts were available, so these norms were not treated as semantic-margin evidence. Familiarity, age of acquisition, log frequency, broad typicality, and word length alone did not significantly predict $\rho$ in the covered subset. The original semantic margin remained predictive after adding each control separately.
